\documentclass[10pt,a4paper]{article}
\usepackage[spanish,es-nodecimaldot]{babel}
\usepackage[utf8]{inputenc}
\usepackage[T1]{fontenc}
\usepackage{lmodern}
\usepackage[top=1.1cm,bottom=1.1cm,left=1.5cm,right=1.5cm]{geometry}
\usepackage{amsmath,amssymb}
\usepackage[table]{xcolor}
\usepackage{booktabs,array}
\usepackage{enumitem}
\usepackage[normalem]{ulem}
\usepackage[most]{tcolorbox}
\usepackage{pgfplots}
\pgfplotsset{compat=1.18}

\setlength{\parindent}{0pt}
\setlength{\parskip}{2pt}
\pagestyle{empty}

\AtBeginDocument{%
  \setlength{\abovedisplayskip}{3pt}%
  \setlength{\belowdisplayskip}{3pt}%
  \setlength{\abovedisplayshortskip}{2pt}%
  \setlength{\belowdisplayshortskip}{2pt}%
}

\begin{document}

\begin{center}
{\large\bfseries\uline{Prueba H (Kruskal Wallis)}}
\end{center}
\vspace{-2pt}

Alternativa no paramétrica de \textbf{Anova.} Comparo que \textbf{k} medias poblacionales (de k muestras) sean significativamente iguales. \textbf{H0:} u1=u2=...=uk

1.Unimos las k muestras y las ordenamos de forma creciente.

2. A cada valor se le indica a que muestra pertenece (I, II,k) y se le asigna un rango (posición). Si dos valores son iguales, se asigna un promedio de sus rangos.

3. Sumamos los rangos de cada muestra (R1, R2, .., Rk)

4. Aplicamos la fórmula H, si \textbf{ni$>$5.} Aceptamos o rechazamos H0
\[
H = \frac{12}{n\cdot(n+1)}\cdot \sum_{i=1}^{k} \frac{\left(R_i\right)^2}{\left(n_i\right)} - 3\cdot(n+1)
\]

Distribución chi cuadrado con \textbf{k-1 grados de libertad}. \textbf{k} el numero de poblaciones. \textbf{n} es la suma de todos los tamaños muestrales

\vspace{2pt}
\hrule
\vspace{4pt}

\begin{center}
{\large\bfseries EXPLICACION \ PRUEBA DE SUMA DE RANGOS - H}
\end{center}
\vspace{-2pt}

\textbf{4.2 \ 2. La Prueba H de Kruskal-Wallis (Tres o más muestras independientes)}

¿Qué ocurre si la investigación involucra tres o más grupos independientes? La prueba U se queda corta. Necesitamos generalizarla para $k \geq 3$ grupos. Esto es la \textbf{Prueba H de Kruskal-Wallis}, que es el equivalente no paramétrico de la prueba ANOVA paramétrica de una vía.

\textbf{Ejercicio 4: Comparación de Métodos contra la Corrosión en Cables}

\textbf{Caso Práctico:} Se evalúa la eficacia de tres tratamientos preventivos anticorrosión midiendo la profundidad máxima de las cavidades (en milésimas de pulgada) en piezas de cable sometidas a un ensayo acelerado:
\begin{itemize}[leftmargin=16pt,itemsep=1pt,topsep=1pt,parsep=0pt]
  \item \textbf{Método A ($n_1 = 6$):} 77, 54, 67, 74, 71, 66
  \item \textbf{Método B ($n_2 = 7$):} 60, 41, 59, 65, 62, 64, 52
  \item \textbf{Método C ($n_3 = 5$):} 49, 52, 69, 47, 56
\end{itemize}

Probar si los tres tratamientos difieren significativamente en su eficacia preventiva utilizando un nivel de significación de $\alpha = 0.05$.

\textbf{Resolución Desarrollada Paso a Paso:}
\begin{enumerate}[leftmargin=18pt,itemsep=2pt,topsep=2pt,parsep=0pt]
  \item \textbf{Planteamiento de Hipótesis:}

  $H_0$: Las poblaciones de los tres tratamientos son idénticas ($\tilde{\mu}_A = \tilde{\mu}_B = \tilde{\mu}_C$)\\
  $H_1$: Al menos uno de los tratamientos difiere en su localización central (distinta eficacia)

  \item \textbf{Ordenamiento de Datos y Asignación de Rangos (Manejo de Empates):} Agrupamos todos los datos ($N = 6+7+5 = 18$ observaciones) y los ordenamos de forma creciente: \hfill {\small\textbf{Ordenamos las K muestras como una sola}}

  \begin{center}
  \vspace{-2pt}
  {\small
  \begin{tabular}{l c c c c c c c c c c c c c c c c c c}
  \toprule
  Valor & 41 & 47 & 49 & 52 & 52 & 54 & 56 & 59 & 60 & 62 & 64 & 65 & 66 & 67 & 69 & 71 & 74 & 77 \\
  Grupo & B & C & C & B & C & A & C & B & B & B & B & B & A & A & C & A & A & A \\
  \midrule
  Rango & 1 & 2 & 3 & \textbf{4.5} & \textbf{4.5} & 6 & 7 & 8 & 9 & 10 & 11 & 12 & 13 & 14 & 15 & 16 & 17 & 18 \\
  \bottomrule
  \end{tabular}
  }
  \vspace{-4pt}
  \end{center}

  \begin{tcolorbox}[colback=red!6, colframe=red!70!black, title={\small\bfseries Atención - Error Frecuente}, boxsep=1pt, left=4pt, right=4pt, top=2pt, bottom=2pt]
  {\small
  \textbf{El tratamiento de empates:} En la tabla observamos que el valor 52 aparece dos veces (posiciones ordinales 4 y 5). En estadística no paramétrica, \textbf{jamás se les asignan rangos arbitrarios} (como 4 al primero y 5 al segundo) porque esto inventaría un orden inexistente. En su lugar, a cada elemento empatado se le asigna el promedio de los rangos que ocuparían conjuntamente:
  \[
  \text{Rango asignado} = \frac{4+5}{2} = 4.5
  \]
  Ambas observaciones reciben el rango 4.5 (el cual se reparte entre los sumatorios de los grupos B y C).
  }
  \end{tcolorbox}

  \item \textbf{Suma de Rangos por Grupo ($R_i$):}
  \begin{align*}
  R_A &= 6 + 13 + 14 + 16 + 17 + 18 = \mathbf{84} \\
  R_B &= 1 + 4.5 + 8 + 9 + 10 + 11 + 12 = \mathbf{55.5} \\
  R_C &= 2 + 3 + 4.5 + 7 + 15 = \mathbf{31.5}
  \end{align*}

  \textit{Control aritmético rápido ($N = 18$):}
  \[
  \sum R = 84 + 55.5 + 31.5 = 171
  \quad \text{y} \quad
  \frac{18(19)}{2} = 171 \quad \checkmark\ \text{(Correcto)}
  \]

  \item \textbf{Cálculo del Estadístico H de Kruskal-Wallis:} Utilizamos la ecuación fundamental de la prueba H:
  \begin{align*}
  H &= \frac{12}{N(N+1)} \left( \sum_{i=1}^{k} \frac{R_i^2}{n_i} \right) - 3(N+1) \\
  H &= \frac{12}{18(19)} \left( \frac{84^2}{6} + \frac{55.5^2}{7} + \frac{31.5^2}{5} \right) - 3(19) \\
  H &= \frac{12}{342} \left( \frac{7056}{6} + \frac{3080.25}{7} + \frac{992.25}{5} \right) - 57 \\
  H &= 0.0350877\,(1176 + 440.036 + 198.45) - 57 \\
  H &= 0.0350877\,(1814.486) - 57 \\
  H &= 63.666 - 57 = \mathbf{6.666} \approx \mathbf{6.66}
  \end{align*}

  \item \textbf{Región Crítica y Decisión (Distribución de Referencia):} Bajo la hipótesis nula, si los tamaños muestrales de cada grupo son mayores a 5 ($n_i > 5$, lo cual se cumple ya que tenemos $n_1 = 6, n_2 = 7, n_3 = 5$), el estadístico H se modela mediante la distribución \textbf{Chi-cuadrado} con $k-1$ grados de libertad (donde $k = 3$ grupos, $gl = 2$). Para $\alpha = 0.05$ y $gl = 2$:
  \[
  \chi^2_{\text{crit}} = \chi^2_{0.05,2} = \mathbf{5.991}
  \]
  La regla de decisión establece que se rechaza la hipótesis nula si $H > 5.991$ (Ver figura 3).

  \begin{center}
  \vspace{-2pt}
  \begin{tikzpicture}
  \begin{axis}[
    width=9.5cm, height=4.6cm,
    title={\small\bfseries Distribución de Referencia $\chi^2$ con 2 g.l. (Kruskal-Wallis) \ (Desgaste de Cables de Cobre)},
    xlabel={\footnotesize Estadístico H de Kruskal-Wallis},
    ylabel={\footnotesize Densidad de Probabilidad},
    xmin=0, xmax=10, ymin=0, ymax=0.52,
    xtick={0,2,4,6,8,10},
    tick label style={font=\scriptsize},
    clip=false
  ]
  \addplot[domain=5.991:10, samples=60, draw=none, fill=red!25] {0.5*exp(-x/2)} \closedcycle;
  \addplot[domain=0.001:10, samples=120, thick, blue!40!black] {0.5*exp(-x/2)};
  \draw[red!70!black, dashed] (axis cs:5.991,0) -- (axis cs:5.991,0.5);
  \fill[orange!80!yellow] (axis cs:6.66,0.018) circle (2.2pt);
  \node[draw=orange!80!black, fill=orange!10, align=center, font=\tiny, anchor=south] at (axis cs:6.4,0.14) {H\_calc = 6.66\\ (Cae en Zona de Rechazo)\\ \textcolor{red!70!black}{(Cola Derecha al 5\%)}};
  \draw[gray] (axis cs:6.5,0.13) -- (axis cs:6.66,0.025);
  \end{axis}
  \end{tikzpicture}
  \vspace{-8pt}

  {\footnotesize Figura 3: Gráfica de la distribución Chi-cuadrado con 2 grados de libertad. El estadístico calculado H (6.66) cae en la zona de rechazo a la derecha del valor crítico 5.991.}
  \end{center}

  \item \textbf{Decisión e Interpretación:} Como $H = 6.66 > \chi^2_{\text{crit}} = 5.991$, se rechaza la hipótesis nula $H_0$. \textbf{Interpretación:} A un nivel de significación del 5\%, existe suficiente evidencia para afirmar que los tres tratamientos contra la corrosión de cables no tienen la misma eficacia preventiva. Al analizar la suma de rangos, el Método A (con el rango más alto, indicando mayor profundidad de cavidad) resulta ser el menos eficaz, mientras que el Método C presenta los mejores resultados preventivos (cavidades menos profundas).
\end{enumerate}

\end{document}
